\documentclass[11pt]{article}

\usepackage[margin=1in]{geometry}
\usepackage{amsmath,amssymb}
\usepackage{enumitem}

\setlength{\parindent}{0pt}
\setlength{\parskip}{0.6em}

\begin{document}

\begin{center}
    {\Large \textbf{MATH 4790: Quiz 2}}\\[0.3em]
    February 6, 2026
\end{center}

\vspace{1em}

\textbf{Name:} \hfill \textbf{Score:} \rule{1.5cm}{0.4pt}/15

\vspace{1em}

\hrule
\vspace{1em}

\textbf{Instructions:} Show all work and justify your answers. Only calculators approved for Exam~FM are permitted. Credit will not be given to just guessing the correct answer.

\vspace{1em}

% -------------------- QUESTION 1 (OLD Q3) --------------------

\noindent\textbf{Question 1.}  

Emma deposits $\$100$ into an account that earns interest at a nominal rate of $12\%$ compounded quarterly. Lucy deposits $\$200$ into an account that earns a constant force of interest of $6\%$. At time $n$, both accounts have the same amount. Find the amount both accounts have at time $n$.

\medskip

\begin{center}

\textbf{(A)} $275.5$
\hspace{1cm}
\textbf{(B)} $300.5$
\hspace{1cm}
\textbf{(C)} $319.5$
\hspace{1cm}
\textbf{(D)} $408.5$
\hspace{1cm}
\textbf{(E)} $467.5$

\end{center}

\newpage

% -------------------- QUESTION 2 --------------------

\textbf{Question 2.}

Two accounts accumulate interest as follows.

\medskip

\textbf{Account 1:} A deposit of \$200 is made at time $t=0$ and the account earns simple interest at $6\%$.

\medskip

\textbf{Account 2:} A deposit of \$100 is made at time $t=0$ and the amount function is
\[
A(t)=100+kt^{2}, \qquad t\ge 0.
\]

\medskip

It is given that the amount of interest earned from time $t=1$ to time $t=3$ is the same for both accounts. Find $k$.

\medskip

\begin{center}

\textbf{(A)} $3$
\hspace{1cm}
\textbf{(B)} $4$
\hspace{1cm}
\textbf{(C)} $5$
\hspace{1cm}
\textbf{(D)} $6$
\hspace{1cm}
\textbf{(E)} $7$

\end{center}

\newpage

% -------------------- QUESTION 3 (OLD Q1) --------------------

\textbf{Question 3.}

Tony deposits \$50 into an account at time $t=0$.

\begin{itemize}
    \item During the first year, the account earns interest at an effective discount rate of $5\%$.
    \item During the next $n$ years, the force of interest is given by
    \[
    \delta_t=\frac{2}{4+2t}.
    \]
    \item During the final year, the account earns interest at an effective interest rate of $5\%$.
\end{itemize}

If the final value of the account is \$100, find $n$.

\medskip

\begin{center}

\textbf{(A)} $1.92$
\hspace{1cm}
\textbf{(B)} $2.43$
\hspace{1cm}
\textbf{(C)} $3.05$
\hspace{1cm}
\textbf{(D)} $3.51$
\hspace{1cm}
\textbf{(E)} $5.17$

\end{center}

\end{document}
